%!TEX root = /Users/markelikalderon/Documents/Git/perceptual_self-consciousness/perceptual_self-consciousness.tex
\chapter{Third Excursus} % (fold)
\label{cha:third_excursus}

\section{Aspects of Perception}
\label{sec:procession_and_reversion_as_aspects_of_perception}

Our scribe, having worked late copying, hurries down the hall, only to catch a glimpse of Proclus, through an open door, performing this evening's theurgic rites. In seeing Proclus, their sense of sight proceeds from their \emph {psuchē}, and transcends it. After all, Proclus is distinct from the scribe and their psyhic activity. Insofar as they are aware that they see what they see, the scribe's sight reverts upon their so seeing. What is more, procession and reversion occur in non-separate time. In proceeding to Proclus who transcends the scribe's psychic activity, their sense of sight, at the same time, reverts upon this activity. The scribe perceives that they see while they see what they see. In proceeding to a transcendent object and reverting upon this activity, each unfolds in tandem over time. The scribe's sense of sight proceeds to Proclus in seeing and reverts upon this activity while seeing. 

There are thus two senses in which procession and reversion are aspects of the scribe's perceptual activity.

In the first sense, procession and reversion are two modes of conscious perception. There is the perceiving of the transcendent objects and there is the awareness of so perceiving and each is an aspect of a single activity. That is what makes talk of \emph{sunaisthēsis} apt. There is a single activity, seeing, with two aspects: an intentional aspect, the apprehension of the transcendent object, and a reflexive aspect, the apprehension of that very activity.

In the second sense, procession and reversion, at least within the sensible realm, share a temporal aspect.

% section procession_and_reversion_as_aspects_of_perception (end)

% chapter third_excursus (end)